\section{Evaluation}
\label{sec:results}
In this section we evaluate \method{} across a number of domains.
In particular, we compare \method{} with two of the most successful
recent data attribution techniques: TRAK \citep{park2023trak} and
EK-FAC \citep{grosse2023studying}; see Appendix~\ref{app:baselines} for the
specifics of these baselines.
Across the board, we find that \method{} provides near-perfect
predictions of how model outputs change when we drop data
(at random) from the training set.

\paragraph{Evaluation metric.}
Recall (from Section~\ref{sec:problem_setup}) that the goal of predictive
data attribution is to predict how a model's output changes as a function of
the model's training data. In order to evaluate the quality of these
predictions,
we adopt the {\em linear datamodeling score} (LDS)
\citep{ilyas2022datamodels,park2023trak}
as our evaluation metric. To compute LDS for a given model output function
$f$ and corresponding data attribution method $\hat{f}$, we follow
the following steps:
\begin{enumerate}
  \item Sample $n$ fixed-sized subsets of the training set,
    which we represent as binary data weights
    $\wvec^{(1)}, \ldots, \wvec^{(n)} \in \{0,1\}^N$,
    where $N$ is the number of total training samples.
    Given a drop-out fraction $p \in [0,1]$, we sample each data weight
    vector $\wvec^{(i)}$ by dropping $pN$ random training samples
    from the training set.

  \item For each data weight vector $\wvec_i$:
    \begin{enumerate}
      \item Compute the {\em true} model output function
        $f(\wvec_i)$ by training a model on the training set with
        data weights $\wvec_i$
        and evaluating the measurement of interest on the trained model.
      \item Compute the {\em predicted} model output function
        $\hat{f}(\wvec_i)$ via the data attribution method $\hat{f}$.
    \end{enumerate}

  \item We compute the LDS as the Spearman correlation between the predicted
    output and the true output over all $n$ data weight vectors, i.e.,
    \begin{align}
      \label{eq:lds}
      \text{LDS} =  \rho\!\left(\left[\hat{f}(\wvec_i)\right]_{i=1}^n,
      \left[\phantom{\hat{f}}\!\!\! f(\wvec_i)\right]_{i=1}^n\right).
    \end{align}
\end{enumerate}

\paragraph{Settings.}
We study scenarios spanning computer vision and
language modeling.
Each scenario comprises a training dataset $S$,
a learning algorithm $\mathcal{A}$,
and a test set $S_T$. Accordingly, each scenario defines $|S_T|$ data
attribution
tasks, where task $i$ is to predict the loss of $\mathcal{A}$ on the
$i$-th sample
in $S_T$. For each data attribution method we consider,
we compute the {\em average} LDS \eqref{eq:lds} across tasks.
\begin{itemize}
  \item {\bf ResNet CIFAR-10 training}: We train ResNet-9~\citep{jordan202494}
    models on subsets of the CIFAR-10 train set~\citep{krizhevsky2009learning},
    and aim to predict cross-entropy loss on CIFAR-10 test samples.
  \item {\bf GPT-2 Wikitext fine-tuning}: We fine-tune 125M
    GPT-2 models~\citep{radford2019language} on subsets of
    Wikitext~\citep{foundation2022english}, and aim to predict
    (language modeling) loss on 50 different test samples.
  \item {\bf Gemma-2B instruction-tuning}: We fine-tune Gemma-2B
    \citep{team2024gemma} with LoRA \citep{hu2021lora} on 
    subsets of three combined instruction tuning datasets
    (Flan V2, DOLLY, OpenAssistant-1
    \citep{longpre2023flan,conover2023free,kopf2024openassistant}), aiming to predict 
    loss on MMLU \citep{hendrycks2021measuring} samples.
\end{itemize}
See Appendix~\ref{app:expdetails} for the exact details of each scenario.


\begin{figure}[bp!]
  \centering
  \usetikzlibrary{positioning}
\usetikzlibrary{matrix}
\pgfplotsset{colormap/Set1}

\definecolor{mycolor1}{RGB}{74,124,182}
\definecolor{mycolor2}{RGB}{142,82,159}
\definecolor{mycolor3}{RGB}{239,134,50}

\begin{tikzpicture}
  \begin{groupplot}[
    group style={group size=3 by 1, horizontal sep=0.5cm},
    height=4.3cm,
    width=5.3cm,
    xlabel={Drop Fraction (\%)},
    ylabel={Correlation},
    title style={yshift=-.15cm},
    ytick={0, 0.25, 0.5, 0.75, 1},
    ymin=-0.1,
    xtick={0,5,10,15,20},
    xticklabels={0\%,5\%,10\%,15\%,20\%},
    grid=both,
    cycle list={
      {mycolor1, mark=square*, thick, error bars/.cd, y dir=both, y explicit, error bar style={color=gray, line cap=round}},
      {mycolor2, mark=*, thick, error bars/.cd, y dir=both, y explicit, error bar style={color=gray, line cap=round}},
      {mycolor3, mark=diamond*, thick, error bars/.cd, y dir=both, y explicit, error bar style={color=gray, line cap=round}},
    },
    legend style={at={(3.4,0.5)}, anchor=west, legend columns=1},
  ]
  
  
  \nextgroupplot[title={ResNet-9 (CIFAR-10)}]
  \addplot+ %
  coordinates {
    (1, 0.234) +- (0.063, 0.061)
    (5, 0.249) +- (0.055, 0.061)
    (10, 0.255) +- (0.102, 0.105)
    (20, 0.308) +- (0.080, 0.072)
  };
  \addlegendentry{EKFAC}
  \addplot+ %
  coordinates {
    (1, 0.347) +- (0.098, 0.088)
    (5, 0.362) +- (0.078, 0.077)
    (10, 0.401) +- (0.110, 0.109)
    (20, 0.393) +- (0.109, 0.096)
  };
  \addlegendentry{TRAK}
  \addplot+ %
  coordinates {
    (1, 0.957) +- (0.008, 0.008)
    (5, 0.922) +- (0.017, 0.017)
    (10, 0.897) +- (0.022, 0.021)
    (20, 0.829) +- (0.021, 0.023)
  };
  \addlegendentry{\method}
  
  \nextgroupplot[title={Gemma-2B (IFT)}, ylabel={}, yticklabels=\empty]
  \addplot+ %
  coordinates {
    (1, 0.036) +- (0.0603, 0.0576)
    (5, 0.046) +- (0.0569, 0.0583)
    (10, -0.038) +- (0.0780, 0.0826)
    (20, 0.036) +- (0.0846, 0.0882)
  };
  \addplot+ %
  coordinates {
    (1, 0.013) +- (0.0716, 0.0798)
    (5, 0.049) +- (0.0745, 0.0754)
    (10, -0.014) +- (0.0799, 0.0812)
    (20, 0.012) +- (0.0856, 0.0890)
  };
  \addplot+ %
  coordinates {
    (1, 0.970) +- (0.0096, 0.0072)
    (5, 0.905) +- (0.0214, 0.0193)
    (10, 0.828) +- (0.0324, 0.0287)
    (20, 0.725) +- (0.0614, 0.0498)
  };
  
  
  
  \nextgroupplot[title={GPT-2 (Wikitext)}, ylabel={}, yticklabels=\empty]
  \addplot+ %
  coordinates {
    (1, 0.342) +- (0.055, 0.057)
    (5, 0.362) +- (0.056, 0.055)
    (10, 0.362) +- (0.051, 0.054)
    (20, 0.362) +- (0.046, 0.045)
  };
  \addplot+ %
  coordinates {
    (1, -0.002) +- (0.029, 0.029)
    (5, 0.030) +- (0.039, 0.039)
    (10, 0.026) +- (0.046, 0.045)
    (20, 0.003) +- (0.037, 0.038)
  };
  \addplot+ %
  coordinates {
    (1, 0.910) +- (0.011, 0.011)
    (5, 0.973) +- (0.003, 0.003)
    (10, 0.974) +- (0.003, 0.003)
    (20, 0.970) +- (0.005, 0.004)
  };
  
  \end{groupplot}
  
  \end{tikzpicture}

  \caption{Linear datamodeling score (LDS) vs. drop fraction across settings for
    \method{} and baselines. The estimates of \method{} consistently correlate
    with the true model outputs (LDS: near $1.0$ for small enough drop fraction)
    while baselines often do not (LDS: below $0.4$). LDS decreases with
    increasing drop fraction for \method{} (as the Taylor estimate moves further
  from the center).}
  \label{fig:cifar_lds}
\end{figure}

\subsection{Results}

As shown in Figure~\ref{fig:cifar_lds}, \method{} attains near-perfect
LDS across settings and drop-out fractions,
although our predictions slightly degrade as we drop more data.
Existing baselines are noisy in comparison; these methods' predictions only
weakly correlate with the ground truth model losses.

In Figure~\ref{fig:scatters}, we randomly select a test example 
from each scenario, and plot predictions of test loss against true test loss for
each data attribution method. \method{} almost exactly predicts the true test
loss, even in absolute terms. On the other hand, baselines barely correlate
with the true predictions, and are mis-scaled in absolute terms
(TRAK and EKFAC predictions are not of the right order of magnitude,
so we rescale them to visualize them on the same plot).

\medskip
\noindent \textbf{Optimality of \method{}.} We observe that the performance of
\method{} degrades with the fraction of samples dropped. While \method{} has
near perfect LDS when predicting the effect of removing a small fraction of the
training data (i.e., 1\%), the LDS degrades at larger drop-out fractions
(i.e., 20\%). Our near-perfect performance when dropping only a few points
indicates that our method is in some sense optimal among linear predictors: 
we can perfectly predict in a small ball around ``not dropping out any points,'' 
but curvature (in training data weight space) 
causes the linear approximation to degrade
further away from training on all the data.

\begin{figure}[ht!]
  \newcommand{\dataset}{cifar}
  \newcommand{\exampleid}{0}
  \newcommand{\scatteropacity}{0.5}
  \newcommand{\setting}{ResNet-9 on CIFAR-10}
  \newcommand{\fracone}{1}
  \newcommand{\frachtwo}{5}
  \newcommand{\datacsvscatterone}{plots/scatter/scatter_results_cifar/ind_0_dropfrac_0.01.csv}
  \newcommand{\datacsvscattertwo}{plots/scatter/scatter_results_cifar/ind_0_dropfrac_0.05.csv}
  \pgfplotsset{colormap/Set1}

\definecolor{Set1Red}{RGB}{74,124,182}
\definecolor{Set1Blue}{RGB}{142,82,159}
\definecolor{Set1Green}{RGB}{239,134,50}

\usetikzlibrary{calc}

%
\def\datacsvcorrelations{plots/scatter/correlations_\dataset/dropfrac_\fracone/example_\exampleid.csv}

%
\pgfplotstableread[col sep=comma]{\datacsvcorrelations}\datatable

%
\pgfplotstablegetelem{0}{corr}\of\datatable
\pgfmathprintnumberto[fixed, precision=2]{\pgfplotsretval}{\trakcorrone}

\pgfplotstablegetelem{1}{corr}\of\datatable
\pgfmathprintnumberto[fixed, precision=2]{\pgfplotsretval}{\ekfaccorrone}

\pgfplotstablegetelem{2}{corr}\of\datatable
\pgfmathprintnumberto[fixed, precision=2]{\pgfplotsretval}{\methodcorrone}

%
\def\dataset{wikitext}
\def\datacsvcorrelations{plots/scatter/correlations_\dataset/dropfrac_\frachtwo/example_\exampleid.csv}

%
\pgfplotstableread[col sep=comma]{\datacsvcorrelations}\datatable

%
\pgfplotstablegetelem{0}{corr}\of\datatable
\pgfmathprintnumberto[fixed, precision=2]{\pgfplotsretval}{\trakcorrhtwo}

\pgfplotstablegetelem{1}{corr}\of\datatable
\pgfmathprintnumberto[fixed, precision=2]{\pgfplotsretval}{\ekfaccorrhtwo}

\pgfplotstablegetelem{2}{corr}\of\datatable
\pgfmathprintnumberto[fixed, precision=2]{\pgfplotsretval}{\methodcorrhtwo}

%


%
\ifdefined\hidelegend
    \pgfplotsset{conditional legend style/.style={legend style={draw=none, fill=none}}}
\else
    \pgfplotsset{conditional legend style/.style={
        legend style={
            at={(0.12,1.22)}, %
            anchor=south, %
            legend columns=4, %
            /tikz/every even column/.append style={column sep=0.5cm}, %
            inner xsep=24pt, %
            inner ysep=2.5pt,  %
            legend image post style={scale=2.5} %
        }
    }}
\fi

\begin{tikzpicture}
    %
    \node (settingbox) [minimum width=0.33\textwidth, 
                        minimum height=0.2\textwidth, 
                        anchor=north west, 
                        align=center] {\setting \\[.5em]
                        Correlation {\bf on Test Example \#$\mathbf{\exampleid}$} \\[.5em]
                        
        \begin{tabular}{ccc}
            \toprule
            Method & $\rho\ (\fracone\%)$ & $\rho\ (\frachtwo\%)$ \\
            \midrule
            EK-FAC & $\ekfaccorrone$ & $\ekfaccorrhtwo$ \\
            TRAK & $\trakcorrone$ & $\trakcorrhtwo$ \\
            \method & $\mathbf{\methodcorrone}$ & $\mathbf{\methodcorrhtwo}$ \\
            \bottomrule
        \end{tabular}
    };

    %
    \begin{axis}[
        name=plot1,
        at={($(settingbox.east)+(2cm,0)$)}, %
        anchor=west, %
        width=0.33\textwidth,
        height=0.33\textwidth,
        xlabel={Predicted Loss},
        ylabel={True Loss},
        title={Drop \fracone\%},
        title style={yshift=-.15cm},
        %
        conditional legend style,
        grid=both,
        grid style={line width=.1pt, draw=gray!10},
        major grid style={line width=.2pt,draw=gray!50},
        scaled y ticks=true,
        yticklabel style={/pgf/number format/.cd, fixed, precision=3},
        xticklabel style={/pgf/number format/.cd, fixed, precision=3},
        after end axis/.code={
            %
            \pgfmathsetmacro{\myxmin}{\pgfkeysvalueof{/pgfplots/xmin}}
            \pgfmathsetmacro{\myxmax}{\pgfkeysvalueof{/pgfplots/xmax}}
            \pgfmathsetmacro{\myymin}{\pgfkeysvalueof{/pgfplots/ymin}}
            \pgfmathsetmacro{\myymax}{\pgfkeysvalueof{/pgfplots/ymax}}
            %
            \pgfmathsetmacro{\linestart}{max(\myxmin, \myymin)}
            \pgfmathsetmacro{\lineend}{min(\myxmax, \myymax)}
            %
            \draw[dashed, gray, very thick, opacity=0.5] 
                (axis cs:\linestart,\linestart) -- 
                (axis cs:\lineend,\lineend);
            }
    ]
        %
        \addplot[
            only marks,
            mark=square*,
            mark size=1.5pt,
            color=Set1Red,
            fill opacity=\scatteropacity,
            draw opacity=\scatteropacity,
            point meta=explicit,
        ] table[x=kron_infl_10, y=true, meta expr=1, col sep=comma] {\datacsvscatterone};
        %
        \ifdefined\hidelegend\else
            \addlegendentry{EK-FAC (re-scaled)}
        \fi
        
        %
        \addplot[
            only marks,
            mark=triangle*,
            mark size=1.5pt,
            color=Set1Blue,
            fill opacity=\scatteropacity,
            draw opacity=\scatteropacity,
            point meta=explicit,
        ] table[x=trak_infl_10, y=true, meta expr=2, col sep=comma] {\datacsvscatterone};
        %
        \ifdefined\hidelegend\else
            \addlegendentry{TRAK (re-scaled)}
        \fi
        
        %
        \addplot[
            only marks,
            mark=*,
            mark size=2pt,
            color=Set1Green,
            fill opacity=1,
            draw opacity=1,
            point meta=explicit,
        ] table[x=infls, y=true, meta expr=0, col sep=comma] {\datacsvscatterone};
        %
        \ifdefined\hidelegend\else
            \addlegendentry{\method{} (unscaled)}
        \fi

        %
        \ifdefined\hidelegend\else
            \addlegendimage{dashed, gray, very thick,
                legend image code/.code={
                    \draw[dashed, gray, very thick] (0cm,0cm) -- (0.32cm,0cm); %
                }
            }
            \addlegendentry{\ $y=x$}
        \fi
    \end{axis}
    
    %
    \begin{axis}[
        name=plot2,
        at={($(plot1.east)+(1.2cm,0)$)},
        anchor=west,
        width=0.33\textwidth,
        height=0.33\textwidth,
        xlabel={Predicted Loss},
        ylabel={},
        title={Drop \frachtwo\%},
        title style={yshift=-.15cm},
        grid=both,
        grid style={line width=.1pt, draw=gray!10},
        major grid style={line width=.2pt,draw=gray!50},
        scaled y ticks=true,
        yticklabel style={/pgf/number format/.cd, fixed, precision=3},
        xticklabel style={/pgf/number format/.cd, fixed, precision=3},
        after end axis/.code={
            %
            \pgfmathsetmacro{\myxmin}{\pgfkeysvalueof{/pgfplots/xmin}}
            \pgfmathsetmacro{\myxmax}{\pgfkeysvalueof{/pgfplots/xmax}}
            \pgfmathsetmacro{\myymin}{\pgfkeysvalueof{/pgfplots/ymin}}
            \pgfmathsetmacro{\myymax}{\pgfkeysvalueof{/pgfplots/ymax}}
            %
            \pgfmathsetmacro{\linestart}{max(\myxmin, \myymin)}
            \pgfmathsetmacro{\lineend}{min(\myxmax, \myymax)}
            %
            \draw[dashed, gray, very thick, opacity=0.5] 
                (axis cs:\linestart,\linestart) -- 
                (axis cs:\lineend,\lineend);
            }
    ]
        %
        \addplot[
            only marks,
            mark=square*,
            mark size=1.5pt,
            color=Set1Red,
            fill opacity=\scatteropacity,
            draw opacity=\scatteropacity,
            point meta=explicit,
        ] table[x=kron_infl_10, y=true, meta expr=1, col sep=comma] {\datacsvscattertwo};
        
        \addplot[
            only marks,
            mark=triangle*,
            mark size=1.5pt,
            color=Set1Blue,
            fill opacity=\scatteropacity,
            draw opacity=\scatteropacity,
            point meta=explicit,
        ] table[x=trak_infl_10, y=true, meta expr=2, col sep=comma] {\datacsvscattertwo};
        
        \addplot[
            only marks,
            mark=*,
            mark size=2pt,
            color=Set1Green,
            fill opacity=1,
            draw opacity=1,
            point meta=explicit,
        ] table[x=infls, y=true, meta expr=0, col sep=comma] {\datacsvscattertwo};
    \end{axis}
    
\end{tikzpicture}

  \renewcommand{\dataset}{gemma}
  \renewcommand{\exampleid}{17}
  \renewcommand{\scatteropacity}{0.5}
  \renewcommand{\setting}{Gemma-2B on IFT data}
  \newcommand{\hidelegend}{true}
  \renewcommand{\fracone}{1}
  \renewcommand{\frachtwo}{5}
  \renewcommand{\datacsvscatterone}{plots/scatter/scatter_results_gemma/example_17_dropfrac_01.csv}
  \renewcommand{\datacsvscattertwo}{plots/scatter/scatter_results_gemma/example_17_dropfrac_05.csv}
  \pgfplotsset{colormap/Set1}

\definecolor{Set1Red}{RGB}{74,124,182}
\definecolor{Set1Blue}{RGB}{142,82,159}
\definecolor{Set1Green}{RGB}{239,134,50}

\usetikzlibrary{calc}

%
\def\datacsvcorrelations{plots/scatter/correlations_\dataset/dropfrac_\fracone/example_\exampleid.csv}

%
\pgfplotstableread[col sep=comma]{\datacsvcorrelations}\datatable

%
\pgfplotstablegetelem{0}{corr}\of\datatable
\pgfmathprintnumberto[fixed, precision=2]{\pgfplotsretval}{\trakcorrone}

\pgfplotstablegetelem{1}{corr}\of\datatable
\pgfmathprintnumberto[fixed, precision=2]{\pgfplotsretval}{\ekfaccorrone}

\pgfplotstablegetelem{2}{corr}\of\datatable
\pgfmathprintnumberto[fixed, precision=2]{\pgfplotsretval}{\methodcorrone}

%
\def\dataset{wikitext}
\def\datacsvcorrelations{plots/scatter/correlations_\dataset/dropfrac_\frachtwo/example_\exampleid.csv}

%
\pgfplotstableread[col sep=comma]{\datacsvcorrelations}\datatable

%
\pgfplotstablegetelem{0}{corr}\of\datatable
\pgfmathprintnumberto[fixed, precision=2]{\pgfplotsretval}{\trakcorrhtwo}

\pgfplotstablegetelem{1}{corr}\of\datatable
\pgfmathprintnumberto[fixed, precision=2]{\pgfplotsretval}{\ekfaccorrhtwo}

\pgfplotstablegetelem{2}{corr}\of\datatable
\pgfmathprintnumberto[fixed, precision=2]{\pgfplotsretval}{\methodcorrhtwo}

%


%
\ifdefined\hidelegend
    \pgfplotsset{conditional legend style/.style={legend style={draw=none, fill=none}}}
\else
    \pgfplotsset{conditional legend style/.style={
        legend style={
            at={(0.12,1.22)}, %
            anchor=south, %
            legend columns=4, %
            /tikz/every even column/.append style={column sep=0.5cm}, %
            inner xsep=24pt, %
            inner ysep=2.5pt,  %
            legend image post style={scale=2.5} %
        }
    }}
\fi

\begin{tikzpicture}
    %
    \node (settingbox) [minimum width=0.33\textwidth, 
                        minimum height=0.2\textwidth, 
                        anchor=north west, 
                        align=center] {\setting \\[.5em]
                        Correlation {\bf on Test Example \#$\mathbf{\exampleid}$} \\[.5em]
                        
        \begin{tabular}{ccc}
            \toprule
            Method & $\rho\ (\fracone\%)$ & $\rho\ (\frachtwo\%)$ \\
            \midrule
            EK-FAC & $\ekfaccorrone$ & $\ekfaccorrhtwo$ \\
            TRAK & $\trakcorrone$ & $\trakcorrhtwo$ \\
            \method & $\mathbf{\methodcorrone}$ & $\mathbf{\methodcorrhtwo}$ \\
            \bottomrule
        \end{tabular}
    };

    %
    \begin{axis}[
        name=plot1,
        at={($(settingbox.east)+(2cm,0)$)}, %
        anchor=west, %
        width=0.33\textwidth,
        height=0.33\textwidth,
        xlabel={Predicted Loss},
        ylabel={True Loss},
        title={Drop \fracone\%},
        title style={yshift=-.15cm},
        %
        conditional legend style,
        grid=both,
        grid style={line width=.1pt, draw=gray!10},
        major grid style={line width=.2pt,draw=gray!50},
        scaled y ticks=true,
        yticklabel style={/pgf/number format/.cd, fixed, precision=3},
        xticklabel style={/pgf/number format/.cd, fixed, precision=3},
        after end axis/.code={
            %
            \pgfmathsetmacro{\myxmin}{\pgfkeysvalueof{/pgfplots/xmin}}
            \pgfmathsetmacro{\myxmax}{\pgfkeysvalueof{/pgfplots/xmax}}
            \pgfmathsetmacro{\myymin}{\pgfkeysvalueof{/pgfplots/ymin}}
            \pgfmathsetmacro{\myymax}{\pgfkeysvalueof{/pgfplots/ymax}}
            %
            \pgfmathsetmacro{\linestart}{max(\myxmin, \myymin)}
            \pgfmathsetmacro{\lineend}{min(\myxmax, \myymax)}
            %
            \draw[dashed, gray, very thick, opacity=0.5] 
                (axis cs:\linestart,\linestart) -- 
                (axis cs:\lineend,\lineend);
            }
    ]
        %
        \addplot[
            only marks,
            mark=square*,
            mark size=1.5pt,
            color=Set1Red,
            fill opacity=\scatteropacity,
            draw opacity=\scatteropacity,
            point meta=explicit,
        ] table[x=kron_infl_10, y=true, meta expr=1, col sep=comma] {\datacsvscatterone};
        %
        \ifdefined\hidelegend\else
            \addlegendentry{EK-FAC (re-scaled)}
        \fi
        
        %
        \addplot[
            only marks,
            mark=triangle*,
            mark size=1.5pt,
            color=Set1Blue,
            fill opacity=\scatteropacity,
            draw opacity=\scatteropacity,
            point meta=explicit,
        ] table[x=trak_infl_10, y=true, meta expr=2, col sep=comma] {\datacsvscatterone};
        %
        \ifdefined\hidelegend\else
            \addlegendentry{TRAK (re-scaled)}
        \fi
        
        %
        \addplot[
            only marks,
            mark=*,
            mark size=2pt,
            color=Set1Green,
            fill opacity=1,
            draw opacity=1,
            point meta=explicit,
        ] table[x=infls, y=true, meta expr=0, col sep=comma] {\datacsvscatterone};
        %
        \ifdefined\hidelegend\else
            \addlegendentry{\method{} (unscaled)}
        \fi

        %
        \ifdefined\hidelegend\else
            \addlegendimage{dashed, gray, very thick,
                legend image code/.code={
                    \draw[dashed, gray, very thick] (0cm,0cm) -- (0.32cm,0cm); %
                }
            }
            \addlegendentry{\ $y=x$}
        \fi
    \end{axis}
    
    %
    \begin{axis}[
        name=plot2,
        at={($(plot1.east)+(1.2cm,0)$)},
        anchor=west,
        width=0.33\textwidth,
        height=0.33\textwidth,
        xlabel={Predicted Loss},
        ylabel={},
        title={Drop \frachtwo\%},
        title style={yshift=-.15cm},
        grid=both,
        grid style={line width=.1pt, draw=gray!10},
        major grid style={line width=.2pt,draw=gray!50},
        scaled y ticks=true,
        yticklabel style={/pgf/number format/.cd, fixed, precision=3},
        xticklabel style={/pgf/number format/.cd, fixed, precision=3},
        after end axis/.code={
            %
            \pgfmathsetmacro{\myxmin}{\pgfkeysvalueof{/pgfplots/xmin}}
            \pgfmathsetmacro{\myxmax}{\pgfkeysvalueof{/pgfplots/xmax}}
            \pgfmathsetmacro{\myymin}{\pgfkeysvalueof{/pgfplots/ymin}}
            \pgfmathsetmacro{\myymax}{\pgfkeysvalueof{/pgfplots/ymax}}
            %
            \pgfmathsetmacro{\linestart}{max(\myxmin, \myymin)}
            \pgfmathsetmacro{\lineend}{min(\myxmax, \myymax)}
            %
            \draw[dashed, gray, very thick, opacity=0.5] 
                (axis cs:\linestart,\linestart) -- 
                (axis cs:\lineend,\lineend);
            }
    ]
        %
        \addplot[
            only marks,
            mark=square*,
            mark size=1.5pt,
            color=Set1Red,
            fill opacity=\scatteropacity,
            draw opacity=\scatteropacity,
            point meta=explicit,
        ] table[x=kron_infl_10, y=true, meta expr=1, col sep=comma] {\datacsvscattertwo};
        
        \addplot[
            only marks,
            mark=triangle*,
            mark size=1.5pt,
            color=Set1Blue,
            fill opacity=\scatteropacity,
            draw opacity=\scatteropacity,
            point meta=explicit,
        ] table[x=trak_infl_10, y=true, meta expr=2, col sep=comma] {\datacsvscattertwo};
        
        \addplot[
            only marks,
            mark=*,
            mark size=2pt,
            color=Set1Green,
            fill opacity=1,
            draw opacity=1,
            point meta=explicit,
        ] table[x=infls, y=true, meta expr=0, col sep=comma] {\datacsvscattertwo};
    \end{axis}
    
\end{tikzpicture}

  \renewcommand{\dataset}{wikitext}
  \renewcommand{\exampleid}{32}
  \renewcommand{\scatteropacity}{0.5}
  \renewcommand{\setting}{GPT-2 on WikiText}
  \renewcommand{\hidelegend}{true}
  \renewcommand{\fracone}{1}
  \renewcommand{\frachtwo}{5}
  \renewcommand{\datacsvscatterone}{plots/scatter/scatter_results_wikitext/example_32_dropfrac_0.01.csv}
  \renewcommand{\datacsvscattertwo}{plots/scatter/scatter_results_wikitext/example_32_dropfrac_0.05.csv}
  \pgfplotsset{colormap/Set1}

\definecolor{Set1Red}{RGB}{74,124,182}
\definecolor{Set1Blue}{RGB}{142,82,159}
\definecolor{Set1Green}{RGB}{239,134,50}

\usetikzlibrary{calc}

%
\def\datacsvcorrelations{plots/scatter/correlations_\dataset/dropfrac_\fracone/example_\exampleid.csv}

%
\pgfplotstableread[col sep=comma]{\datacsvcorrelations}\datatable

%
\pgfplotstablegetelem{0}{corr}\of\datatable
\pgfmathprintnumberto[fixed, precision=2]{\pgfplotsretval}{\trakcorrone}

\pgfplotstablegetelem{1}{corr}\of\datatable
\pgfmathprintnumberto[fixed, precision=2]{\pgfplotsretval}{\ekfaccorrone}

\pgfplotstablegetelem{2}{corr}\of\datatable
\pgfmathprintnumberto[fixed, precision=2]{\pgfplotsretval}{\methodcorrone}

%
\def\dataset{wikitext}
\def\datacsvcorrelations{plots/scatter/correlations_\dataset/dropfrac_\frachtwo/example_\exampleid.csv}

%
\pgfplotstableread[col sep=comma]{\datacsvcorrelations}\datatable

%
\pgfplotstablegetelem{0}{corr}\of\datatable
\pgfmathprintnumberto[fixed, precision=2]{\pgfplotsretval}{\trakcorrhtwo}

\pgfplotstablegetelem{1}{corr}\of\datatable
\pgfmathprintnumberto[fixed, precision=2]{\pgfplotsretval}{\ekfaccorrhtwo}

\pgfplotstablegetelem{2}{corr}\of\datatable
\pgfmathprintnumberto[fixed, precision=2]{\pgfplotsretval}{\methodcorrhtwo}

%


%
\ifdefined\hidelegend
    \pgfplotsset{conditional legend style/.style={legend style={draw=none, fill=none}}}
\else
    \pgfplotsset{conditional legend style/.style={
        legend style={
            at={(0.12,1.22)}, %
            anchor=south, %
            legend columns=4, %
            /tikz/every even column/.append style={column sep=0.5cm}, %
            inner xsep=24pt, %
            inner ysep=2.5pt,  %
            legend image post style={scale=2.5} %
        }
    }}
\fi

\begin{tikzpicture}
    %
    \node (settingbox) [minimum width=0.33\textwidth, 
                        minimum height=0.2\textwidth, 
                        anchor=north west, 
                        align=center] {\setting \\[.5em]
                        Correlation {\bf on Test Example \#$\mathbf{\exampleid}$} \\[.5em]
                        
        \begin{tabular}{ccc}
            \toprule
            Method & $\rho\ (\fracone\%)$ & $\rho\ (\frachtwo\%)$ \\
            \midrule
            EK-FAC & $\ekfaccorrone$ & $\ekfaccorrhtwo$ \\
            TRAK & $\trakcorrone$ & $\trakcorrhtwo$ \\
            \method & $\mathbf{\methodcorrone}$ & $\mathbf{\methodcorrhtwo}$ \\
            \bottomrule
        \end{tabular}
    };

    %
    \begin{axis}[
        name=plot1,
        at={($(settingbox.east)+(2cm,0)$)}, %
        anchor=west, %
        width=0.33\textwidth,
        height=0.33\textwidth,
        xlabel={Predicted Loss},
        ylabel={True Loss},
        title={Drop \fracone\%},
        title style={yshift=-.15cm},
        %
        conditional legend style,
        grid=both,
        grid style={line width=.1pt, draw=gray!10},
        major grid style={line width=.2pt,draw=gray!50},
        scaled y ticks=true,
        yticklabel style={/pgf/number format/.cd, fixed, precision=3},
        xticklabel style={/pgf/number format/.cd, fixed, precision=3},
        after end axis/.code={
            %
            \pgfmathsetmacro{\myxmin}{\pgfkeysvalueof{/pgfplots/xmin}}
            \pgfmathsetmacro{\myxmax}{\pgfkeysvalueof{/pgfplots/xmax}}
            \pgfmathsetmacro{\myymin}{\pgfkeysvalueof{/pgfplots/ymin}}
            \pgfmathsetmacro{\myymax}{\pgfkeysvalueof{/pgfplots/ymax}}
            %
            \pgfmathsetmacro{\linestart}{max(\myxmin, \myymin)}
            \pgfmathsetmacro{\lineend}{min(\myxmax, \myymax)}
            %
            \draw[dashed, gray, very thick, opacity=0.5] 
                (axis cs:\linestart,\linestart) -- 
                (axis cs:\lineend,\lineend);
            }
    ]
        %
        \addplot[
            only marks,
            mark=square*,
            mark size=1.5pt,
            color=Set1Red,
            fill opacity=\scatteropacity,
            draw opacity=\scatteropacity,
            point meta=explicit,
        ] table[x=kron_infl_10, y=true, meta expr=1, col sep=comma] {\datacsvscatterone};
        %
        \ifdefined\hidelegend\else
            \addlegendentry{EK-FAC (re-scaled)}
        \fi
        
        %
        \addplot[
            only marks,
            mark=triangle*,
            mark size=1.5pt,
            color=Set1Blue,
            fill opacity=\scatteropacity,
            draw opacity=\scatteropacity,
            point meta=explicit,
        ] table[x=trak_infl_10, y=true, meta expr=2, col sep=comma] {\datacsvscatterone};
        %
        \ifdefined\hidelegend\else
            \addlegendentry{TRAK (re-scaled)}
        \fi
        
        %
        \addplot[
            only marks,
            mark=*,
            mark size=2pt,
            color=Set1Green,
            fill opacity=1,
            draw opacity=1,
            point meta=explicit,
        ] table[x=infls, y=true, meta expr=0, col sep=comma] {\datacsvscatterone};
        %
        \ifdefined\hidelegend\else
            \addlegendentry{\method{} (unscaled)}
        \fi

        %
        \ifdefined\hidelegend\else
            \addlegendimage{dashed, gray, very thick,
                legend image code/.code={
                    \draw[dashed, gray, very thick] (0cm,0cm) -- (0.32cm,0cm); %
                }
            }
            \addlegendentry{\ $y=x$}
        \fi
    \end{axis}
    
    %
    \begin{axis}[
        name=plot2,
        at={($(plot1.east)+(1.2cm,0)$)},
        anchor=west,
        width=0.33\textwidth,
        height=0.33\textwidth,
        xlabel={Predicted Loss},
        ylabel={},
        title={Drop \frachtwo\%},
        title style={yshift=-.15cm},
        grid=both,
        grid style={line width=.1pt, draw=gray!10},
        major grid style={line width=.2pt,draw=gray!50},
        scaled y ticks=true,
        yticklabel style={/pgf/number format/.cd, fixed, precision=3},
        xticklabel style={/pgf/number format/.cd, fixed, precision=3},
        after end axis/.code={
            %
            \pgfmathsetmacro{\myxmin}{\pgfkeysvalueof{/pgfplots/xmin}}
            \pgfmathsetmacro{\myxmax}{\pgfkeysvalueof{/pgfplots/xmax}}
            \pgfmathsetmacro{\myymin}{\pgfkeysvalueof{/pgfplots/ymin}}
            \pgfmathsetmacro{\myymax}{\pgfkeysvalueof{/pgfplots/ymax}}
            %
            \pgfmathsetmacro{\linestart}{max(\myxmin, \myymin)}
            \pgfmathsetmacro{\lineend}{min(\myxmax, \myymax)}
            %
            \draw[dashed, gray, very thick, opacity=0.5] 
                (axis cs:\linestart,\linestart) -- 
                (axis cs:\lineend,\lineend);
            }
    ]
        %
        \addplot[
            only marks,
            mark=square*,
            mark size=1.5pt,
            color=Set1Red,
            fill opacity=\scatteropacity,
            draw opacity=\scatteropacity,
            point meta=explicit,
        ] table[x=kron_infl_10, y=true, meta expr=1, col sep=comma] {\datacsvscattertwo};
        
        \addplot[
            only marks,
            mark=triangle*,
            mark size=1.5pt,
            color=Set1Blue,
            fill opacity=\scatteropacity,
            draw opacity=\scatteropacity,
            point meta=explicit,
        ] table[x=trak_infl_10, y=true, meta expr=2, col sep=comma] {\datacsvscattertwo};
        
        \addplot[
            only marks,
            mark=*,
            mark size=2pt,
            color=Set1Green,
            fill opacity=1,
            draw opacity=1,
            point meta=explicit,
        ] table[x=infls, y=true, meta expr=0, col sep=comma] {\datacsvscattertwo};
    \end{axis}
    
\end{tikzpicture}

  \caption{Results of \method{} and baselines on randomly chosen, individual
    samples from the three settings we consider: CIFAR-10, Gemma-2B, and GPT-2.
    We evaluate by predicting model output after dropping a random 1\%/5\% of
    the data (cf. \eqref{eq:lds}) and plotting the results against the true
    model output for that drop set. \method{} estimates consistently highly
    correlate with the true output across settings (for small enough training
  data drop fractions).}
  \label{fig:scatters}
\end{figure}

